% Oberdieck--Pixton primitivity scope.
% Primitive K3 classes, imprimitive residues, and Conjecture B.

\providecommand{\kBKM}{\kappa_{\mathrm{BKM}}}
\providecommand{\kch}{\kappa_{\mathrm{ch}}}
\providecommand{\kcat}{\kappa_{\mathrm{cat}}}
\providecommand{\kfib}{\kappa_{\mathrm{fiber}}}
\providecommand{\Zred}{Z^{\mathrm{red}}}
\providecommand{\cN}{\mathcal{N}}
\providecommand{\ZZ}{\mathbb{Z}}
\providecommand{\QQ}{\mathbb{Q}}
\providecommand{\fg}{\mathfrak{g}}
\providecommand{\Phiten}{\Phi_{10}}
\section*{Oberdieck--Pixton Primitivity Scope}
\label{sec:op-primitive-scope}

\subsection*{Primitive $\beta$}

$\beta \in H_2(K3,\ZZ)$ is \emph{primitive} if $\beta = m\beta_0$ with
$m \ge 2$, $\beta_0$ integral, has no solution; equivalently $\beta$
generates $\QQ\beta\cap H_2(K3,\ZZ)$ and is indivisible in
$\Lambda_{K3}\cong U^{\oplus 3}\oplus E_8(-1)^{\oplus 2}$.  Imprimitive
$\beta = m\beta_0$, $m = \mathrm{div}(\beta) \ge 2$.

\subsection*{Geometric Interpretation and the $E$-Factor}

For the Oberdieck--Pixton theorem $\Zred^{GW}(K3\times E) = -1/\Phiten$,
the primitivity hypothesis applies to the \emph{K3-factor class alone};
the $E$-factor degree $d \in H_2(E,\ZZ)\cong\ZZ$ is free, and every
$d \ge 0$ is summed.  Geometrically, a primitive $\beta$ represents a
K3 fibre cycle realisable as a single (not multi-covered) rational nodal
curve in the Beauville--Mukai rank-one Lagrangian fibration; the
Kawai--Yoshioka 2000 count $\sum e(\mathrm{Hilb}^n(K3_\beta))$ that
enters Yau--Zaslow / MPT rests on primitive $\beta$.  Imprimitive
$\beta = m\beta_0$ is a multiple cover of the primitive cycle; its
reduced GW contribution is a sum over divisors of $m$ weighted by
Möbius-type factors from multi-cover reduction (Gopakumar--Vafa form).

The $E$-factor carries no analogous primitivity hypothesis: OP 2017
p.~2 treats the full $(\beta, d)$-sum with $\beta$ primitive in K3 but
$d$ arbitrary, and the Igusa $\Phiten$ encodes the $(p,q,\tilde q)$
triple where $p$ is the $E$-translation, $q$ the K3 arithmetic
genus variable $q = e^{2\pi i\tau}$, $\tilde q$ the $E$-degree variable
$\tilde q = e^{2\pi i\sigma}$.  Primitivity is imposed on the $q$-variable
direction only.

\subsection*{Imprimitive Classes and OP 2014 Conjecture B}

Oberdieck--Pandharipande 2014 (arXiv:1411.1514) Conjecture B states
the all-classes reduction rule: for $\beta = m\beta_0$ with
$\beta_0$ primitive,
\[
  \cN^{X,\mathrm{red}}_{g,\beta,d}
  \;=\; \sum_{k \mid m} k^{2g-3} \cdot f(m/k) \cdot
  \cN^{X,\mathrm{red}}_{g,\beta_0/k^{?},d}
\]
(schematic; the exact form is the multiple-cover formula of
Pandharipande--Thomas / Maulik--Pandharipande--Thomas applied after
reduction).  Partial results:
\begin{itemize}
\item \textbf{MPT 2010} (arXiv:1001.2719): reduces the obstruction
theory on K3 by the $(2,0)$-form, proves the reduced theory is
well-defined on \emph{all} classes $\beta$, primitive or not; the
multiple-cover formula is conjectural at the reduced level.
\item \textbf{PT 2014} (Pandharipande--Thomas, KKV formula,
arXiv:1206.5490): proves the \emph{Katz--Klemm--Vafa} formula
expressing imprimitive K3 GV invariants $n_g(\beta)$ in terms of
primitive ones via a Hecke-like sum; restricted to K3, not
$K3\times E$.
\item \textbf{Bryan--Oberdieck 2018} (specific imprimitive cases):
handles $\beta = 2\beta_0$ at $g=0$ for select K3 lattices via
direct Noether--Lefschetz degeneration; does not extend to all
$m$ or all $g$.
\end{itemize}
The $m \ge 2$ extension of the Igusa cusp form identity remains the
\emph{singular open conjecture} in the OP corpus; OP 2017 p.~5 explicitly
defers it.

\subsection*{Consequence for Conjecture 1'}

Conjecture 1' asserts the three-way equivalence
\[
  \phi_{g_N,h_M}
  \xrightarrow{\text{Borcherds}} \Delta_{k_N}^{g_N,h_M}
  \xrightarrow{\text{MNOP}} Z^{\mathrm{DT}}_{K3\times E}(g_N,h_M)^{-1/2}
  \xrightarrow{\text{denom}} \fg^{\mathrm{BKM}}_{N,h_M}
\]
on the programme core chamber.  The middle MNOP arrow
consumes the Oberdieck--Pixton identity $\Zred^{DT}(K3\times E) =
-1/\Phiten$, which OP 2017 proves \emph{for $\beta$ primitive in
$H_2(K3,\ZZ)$ only}.  The programme-level consequence splits into two
identities:
\begin{enumerate}
\item \emph{Primitive theorem.}  The
root-multiplicity formula $\mathrm{mult}_{\mathrm{BKM}}(\alpha) =
c_{K3}(4nm-\ell^2,\ell)$ at a root $\alpha = (n,\ell,m) \in
\Lambda^{2,1}_{II}$ with K3-projection $(n,\ell)$ corresponding to a
primitive K3 cycle holds via the Borcherds--Gritsenko denominator
identity and the OP 2017 theorem.
\item \emph{Imprimitive conjecture.}  For
$\alpha = (mn_0, m\ell_0, m^2 m_0)$ with $\gcd(n_0,\ell_0) = 1$ and
$m \ge 2$, the analogous identity is
\[
  \mathrm{mult}_{\mathrm{BKM}}(\alpha)
  \;=\; \sum_{k \mid m} \mu\!\big(m/k\big)\,
        c_{K3}\big(4 n_0 m_0 k^2 - \ell_0^2 k^2,\, k\ell_0\big)
        \cdot (\text{correction}),
\]
a Möbius-summed expression whose precise form is OP 2014
Conjecture B.  Neither side is a single Jacobi coefficient; the
identity is a Hecke-type equality between the imprimitive BKM
multiplicity and an arithmetic average of primitive K3 data.
\end{enumerate}

\subsection*{Quantitative Scope}

Charges on $K3\times E$ live in $\Gamma_{3,2} := \Lambda_{K3}\oplus U
\cong U^{\oplus 3}\oplus E_8(-1)^{\oplus 2}\oplus U$ of signature
$(3,2)$ and rank $24$.  The K3-primitive sublocus is
\[
  \Gamma_{3,2}^{\mathrm{prim}}
  \;:=\; \{\gamma = (r,\beta,n) : \beta \in H_2(K3,\ZZ)\
  \text{primitive}\},
\]
which has natural density $6/\pi^2 \approx 0.608$ in rank-$22$ K3
classes (asymptotic density of primitive vectors in a rank-$r$
lattice: $1/\zeta(r) \to 1$; at rank $22$ this is essentially $1$).
The Heegner divisors of charge
discriminant $\Delta = 4nm - \ell^2$ small and negative populate
low-rank sublattices where the primitive fraction can be much smaller.
Concretely:
\begin{itemize}
\item At discriminant $\Delta = -1$ (the Heegner point controlling
the $\Phiten$ lowest polar expansion, i.e.\ the $\tau(a) = 9$
multiplicity), the unique primitive class is $(n,\ell,m) = (1,1,1)$;
all its multiples $(m,m,m)$ for $m \ge 2$ are imprimitive.  The
Conjecture~1' identity at $\Delta = -4m^2 + m^2 = -3m^2$ $(m\ge 2)$
is untreated by OP 2017.
\item At discriminant $\Delta = 0$ (null primitives), OP 2017 covers
the single-fibre multiplicity $c(0) = 20$; imprimitive nulls
$(2n_0,0,2m_0)$ etc.\ require Conjecture B.
\item Across the six-routes meta-diagram for $G(K3\times E)$, routes
(1) Mukai lattice, (3) BKM Borcherds weight, (5) twined elliptic
genus are primitivity-agnostic; routes (2) Igusa $\Phiten$ via MNOP,
(4) K3 fibre-rank via OP-2017 reduced DT, and (6) $\Phi$-functor
vacuum character via MNOP/reduced PT all pass through OP 2017 and
hence inherit the primitive-$\beta$ restriction.
\end{itemize}

\subsection*{Consequences of Conjecture B}

\begin{enumerate}
\item The \emph{Hecke-operator structure} on the BKM denominator:
the expected action of $T_m$ on $\Delta_5$ producing the imprimitive
multiplicity via a Möbius sum is the content of Conjecture B; without
it, the BKM algebra $\fg^{\mathrm{BKM}}_{\Delta_5}$ is known root-by-root
on primitive roots only, not as a Hecke-equivariant structure.
\item The \emph{Conjecture 1' completeness claim on all lattice
charges}: the three-way equivalence holds with a scope
qualifier ``primitive K3 component'' attached to the middle MNOP
arrow; the imprimitive charges remain a conjectural extension.  The
programme's Theorem CY-D universality ($\kch(A_X) = \chi(\mathcal{O}_X)$
on compact CY$_d$) is unaffected (rank-independent), as is
$\kBKM(\Phi_N) = c_N(0)/2$ (Borcherds-weight universality, OP-2017
independent).
\item The \emph{$g$-twined DT generating functions} at $N = 2,3,4,6$:
BOP 2019 and Bryan--Steinberg 2014 extend OP 2017 to
$(g_N,\mathrm{id})$-twined DT at primitive $\beta$; twined + imprimitive
is doubly conjectural.
\end{enumerate}

\subsection*{Scope}

Conjecture 1' holds as stated on the primitive K3-component sublocus
of $\Gamma_{3,2}$, which covers the low-discriminant programme core
$\Delta \in \{-1, 0, \text{small}\}$ fully.  The imprimitive residue,
where $\beta = m\beta_0$ with $m \ge 2$, is the content of OP 2014
Conjecture B and is the single most important open extension blocking
a fully Hecke-equivariant BKM multiplicity formula across all roots.
Every Vol III invocation of $\Zred^{DT}(K3\times E) = -1/\Phiten$
carries the qualifier ``$\beta$ primitive in $H_2(K3,\ZZ)$''. The
imprimitive extension is OP 2014 Conjecture B.

\paragraph{Primary sources.}
Oberdieck--Pixton 2017 arXiv:1706.10100, Invent.\ Math.\ \textbf{222}
(2020) (primitive-$\beta$ Igusa cusp form theorem);
Oberdieck--Pandharipande 2014 arXiv:1411.1514 (Conjectures A--G, with
B the imprimitive reduction);
Maulik--Pandharipande--Thomas 2010 Geom.\ Topol.\ \textbf{14}
arXiv:1001.2719 (reduced obstruction theory on K3, all $\beta$);
Pandharipande--Thomas 2014 arXiv:1206.5490 (KKV formula, K3-only
imprimitive);
Bryan--Oberdieck 2018 (specific imprimitive cases on $K3\times E$);
Gritsenko--Nikulin 1998 Int.\ J.\ Math.\ \textbf{9} (Siegel--Borcherds
lift, primitivity-agnostic on the automorphic side);
Kawai--Yoshioka 2000 (Hilbert-scheme Euler count, primitive $\beta$).
